\chapter*{ЗАКЛЮЧЕНИЕ}
\addcontentsline{toc}{chapter}{ЗАКЛЮЧЕНИЕ}
В ходе работы формализована задача рекомендательных алгоритмов и реализована система рекомендаций рецептов блюд на основе предпочтений и возможностей пользователя.
Также проведено исследование зависимости её характеристик от выбора алгоритма рекомендаций.
В результате исследования выбран алгоритм, при которой вектора рецептов преобразуются в $5$ параметров с помощью профиля пользователя и подаются на вход регрессионной модели, если оценок меньше $5$ и модели <<$k$ ближайших соседей>>, если их не менее $5$.
При параметре регуляризации регрессии $11$ алгоритм позволяет в $80\%$ случаев не рекомендовать пользователю того, что ему не нравится.

Выполнены следующие задачи:
\begin{itemize}
    \item проанализирована предметная область алгоритмов рекомендаций;
    \item  описаны основные методы формирования рекомендаций и выполнить их классификацию;
    \item  проанализированы существующие программы для рекомендаций рецептов;
    \item  выбран метод рекомендаций рецептов;
    \item  спроектирована рекомендательная систему в виде программного обеспечения;
    \item  реализовано программное обеспечение;
    \item реализованы выбранные алгоритмы рекомендаций;
    \item  подготовлен набор данных для обучения и тестирования моделей рекомендаций;
    \item  сравнена точность реализованных моделей рекомендаций на подготовленном наборе данных.
\end{itemize}
Таким образом, задачи работы выполнены и её цель достигнута.